%\newpage
\begin{figure}[b]
\centering
\includegraphics[width=0.7\textwidth]{eh_april91_vcs}
\caption[Measured benzene, toluene and sulphate concentrations (vertical cross-section)
at the Eden Hill field site, after \cite{davis99}.]{Measured benzene, toluene and sulphate concentrations (vertical cross-section) at the Eden Hill field site, after \cite{davis99}.}
\label{fig:eh_april91_vcs}
\end{figure}

%\begin{figure}[tb]
%\centering
%\includegraphics[width=\textwidth]{edenhill}
%\caption[Measured benzene, toluene and sulphate concentrations at the Eden Hill field site %after \cite{davis99}]{Measured benzene, toluene and sulphate concentrations at the Eden Hill %field site after \cite{davis99}}
%\label{fig:edenhill}
%\end{figure}

\section[PHT3D Exercise 4: Microbially mediated petroleum hydrocarbon degradation under sulfate-reducing conditions]{PHT3D Exercise 4: Microbially mediated petroleum hydrocarbon degradation under sulfate-reducing conditions}

Numerical modelling is widely used to analyse and predict the risk associated with
using natural attenuation as a remediation scheme. With few exceptions, most
practical model applications to contaminated field sites involve a vertically averaged two-dimensional simulation model. This assumes that vertical concentration gradients 
over the aquifer depth are negligible. 
However, as pointed out, for example, by \citet{grathwohl2000}, 
transversal vertical mixing can
be a critical (physical) factor for controlling the length of a naturally
attenuating contaminant plume.

Unfortunately, most contaminated sites are not well characterised with respect
to the vertical distribution of dissolved contaminants and inorganic groundwater
constituents within the aquifer. Therefore the identification and quantification
of the vertical mixing process (and how to model it) is generally rather
difficult.

The present exercise is loosely based on the field  \citep[e.g.,][]{barber91,thierrin92,thierrin95,davis99} and corresponding 
modelling studies \citep[e.g.,][]{prommer_tueb, prommer_og, prommer_jch}
that provided a detailed investigation on the fate of hydrocarbon compounds at a 
contaminated site in Perth, Western Australia. 
At this site detailed information of hydrochemical parameters have been intensively
recorded. The data show that toluene, ethylbenzene and xylenes were mineralised under
sulfate-reducing conditions in a seasonally varying groundwater flow field while
benzene was shown to be persistent. Figure \ref{fig:eh_april91_vcs} shows 
measured concentrations of benzene, toluene and
sulphate for a vertical cross-section along the contaminant plume.
The observed long, thin plumes within the relatively
homogeneous sand aquifer indicate low transversal dispersivities, providing a significant
challenge to achieve an accurate numerical description. 

In the exercise you will prepare a simple vertical cross-sectional model 
along a flow path of the contaminant plume within the unconfined aquifer. 
The main purpose of the exercise is to demonstrate the

\begin{itemize}

\item incorporation of microbial kinetics into reactive transport models

\item simulation of multi-component NAPL (Non Aqueous Phase Liquid) dissolution with the PHREEQC/PHT3D framework and the effect of an ageing NAPL source 

\item numerical problems that may be associated with advection dominated reactive transport

\end{itemize}

The exercise has been designed such that all the major phenomena that 
characterise the behaviour of petroleum hydrocarbon plumes occur within
the relatively short simulation period. In reality NAPL concentrations near the water table
might be higher and the mixture of compounds within the NAPL phase can 
also differ from the one used in the present example.

\begin{table}
\caption{Flow and transport parameters used for PHT3D Exercise 4.}
\begin{center}
\begin{tabular}{lc} \hline
Flow simulation                     & steady state \\
Total simulation time $(days)$       & 600 \\
Stress periods                       & 1 \\
Time steps           & 120 \\
Model \  length\     $ (m)                $         &  150            \\
Model \  thickness\  $ (m)                $         &  4            \\
Aquifer top\               $(m)        $         &  20 m          \\
Aquifer bottom\            $(m)        $         &  16 m          \\
Grid \ spacing\      $\Delta x $ (m)                &    5            \\
Grid \ spacing\      $\Delta z $ (m)                &    0.20         \\
Porosity\            $            $         &  0.2           \\
Horizontal hydraulic conductivity\ $(m/day)           $         & 20         \\ 
Vertical   hydraulic conductivity\ $ (m/day)           $         & 20       \\ 
Flux at upstream boundary\ $(m/day)          $         &  0.05          \\ 
Groundwater recharge\ $(m/day)          $         &  0.001          \\ 
Piezometric head downstream boundary\ $(m)        $         &  19 m          \\ 
Longitudinal dispersivity\ $(m)         $         &  0.1          \\  
Transversal  dispersivity\ $(m)         $         &  0.01          \\ \hline

\end{tabular}
\label{edenhill_flow}
\end{center}
\end{table}

\subsection[Setting up the flow problem]{Setting up the flow problem}

A cross-sectional model will be set up.
Groundwater flow in the model domain is driven by both, inflow from the upstream boundary
and by groundwater recharge. To set up the flow model proceed as follows:

\begin{itemize} 

\item Define the model grid according to the discretisation suggested in Table \ref{edenhill_flow}.

\item Set the time units to {\color{red} \textbf{days}} and define a 600 day long stress period as total simulation time, subdivided into 120 {\color{red} \textbf{Time Steps}}. 

\item Define the horizontal/vertical hydraulic conductivity (20 $m/day$) and the porosity (0.2).

\item Define a fixed head boundary condition at the downstream end. Note, a fixed head
boundary condition will not work for grid cells that are positioned above the water table.

\item Allocate suitable initial hydraulic heads.

\item Enter a recharge rate of 0.001 $m/day$ and define a fixed flow boundary with 
a flow rate of 0.05 $m/day$ at the upstream boundary (or 0.01 $m^3/day$ per grid cell).
Also make sure that the recharge is always applied to the highest active grid cell.

\item To prevent drying grid cells activate MODFLOWs (re)wetting capabilities by selecting Models $\rightarrow$ MODFLOW $\rightarrow$ Flow packages $\rightarrow$ Wetting Capability.
Set the wetting factor {\color{red} \textbf{WETFCT}} to a value of 0.1 and select the option 
{\color{red} \textbf{h = BOT + WETFCT $\times \mid$ THRESH $\mid$}}. Set the wetting treshold
{\color{red} \textbf{$\mid$ TRESH $\mid$}} to a value of -.01.

\item As a check your simulated head at the upstream boundary should be 19.56 m.
\end{itemize} 

\subsection[Simulating nonreactive transport]{Simulating nonreactive transport}

Before the reactive transport simulations are started we will briefly check
the setup of the model with a tracer transport simulation. On the first sight 
this step might appear as a waste of time. However, experience shows
that when reactive transport simulations fail to work it is more often than not
caused by a problem related to the flow or the general model setup. 
As a brief check perform the following steps (see also PHT3D Exercise 2 for comparison):

\begin{itemize} 

\item Select Models $\rightarrow$ MT3DMS $\rightarrow$ Simulation Settings and define
a species called {\color{red} \textbf{Tracer}}.

\item Define the settings for advective and dispersive transport. 

\item Set the initial concentrations to 0.

\item Decide what type of boundary conditions would suit the problem.

\item Add additional {\color{red} \textbf{Output Times}} to see intermediate results.

\item Define a contaminant source by defining a concentration value of 1 for the recharge water in the zone between x = $15 m$ and x = $35 m$.

\item Run the simulations and evaluate the plausability of the results. 

\end{itemize}

\subsection[Simulating the dissolution from a multi-component NAPL source]{Simulating the dissolution from a multi-component NAPL source}
 
The first part of the simulation considers only the mass transfer 
process in a NAPL-containing source zone and the 
subsequent advective-dispersive transport. In the source zone, 
dissolution of NAPL compounds,
\citep[e.g.,][]{miller_trapped_90} acts as a continuous contamination source
for the passing groundwater. The rate at which mass is transferred
from the NAPL into the aqueous phase is a function of

\begin{itemize}
    \item Interfacial area between the NAPL phase and the aqueous phase
    \item Extent and morphology of the source (in particular the
maximum cross-sectional area perpendicular to the main groundwater
flow direction)
    \item Groundwater flow velocity
    \item Solubility of individual hydrocarbon compounds
    \item Composition of the NAPL source (mole fraction)
\end{itemize}

Notwithstanding the relative importance of each of these factors,
the concentration of individual hydrocarbon compounds in the
groundwater will, within the contamination source zone, very often
reach an equilibrium concentration that is equal or close to the
multi-component solubility of the compound \citep{eberhardt02}.
This multicomponent solubility $C_{i_{sat},mc}$ is described
by Raoult's law \citep{schwarzb}:
\begin{equation}
{C_{i}^{sat,mc} = C_{i}^{sat} \gamma_{i} m_{i} \ }
\label{eq:raoult2}
\end{equation}
where $C_{i}^{sat}$ is the single-species aqueous-phase
solubility
\citep[available from tabulated sources, e.g.,][]{schwarzb} 
of the organic compound in question within a
mixture of compounds with different physico-chemical properties,
$\gamma_{i}$ is the activity coefficient of the $i^{th}$ organic
compound (typically assumed to be unity) and $m_{i}$ is the
mole fraction of the $i^{th}$ organic compound within the NAPL
mixture. It is defined as
\begin{equation}
{m_i  = \frac {C_n^i} {C_n^{tot}}} \label{eq:raoult_mc}
\end{equation}
where $C_n^i$ is the molar concentration of compound $i$ in the
NAPL phase and $C_n^{tot}$ is the total molar concentration of all
organic compounds in the NAPL phase. A high groundwater flow
velocity, among other factors, might result in a kinetically
limited dissolution of NAPL compounds. The simplest model that
describes the concentration change of the $i^{th}$ compound in
groundwater is
\begin{equation}
{r_{dis_i} = \omega_i (C_{sat,mc}^i - C_i)}
\label{eq:raoult_dis}
\end{equation}
where $C_{i}$ is the concentration of the $i^{th}$ organic
compound in the groundwater and $\omega_i$ is a mass-transfer rate
coefficient that is a product of a mass transfer coefficient and
the specific interfacial area between NAPL phase and water. The
combination of (\ref{eq:raoult2}) and (\ref{eq:raoult_dis})
applies to arbitrary dissolution rates. Note that $\omega$
approaches infinity for equilibrium dissolution. In that case
$C_{i}$ equals $C_{i}^{sat,mc}$.
\par For the modelling exercise we use Eqn. (\ref{eq:raoult_dis}) to
compute NAPL dissolution. The equation was translated into a rate expression 
for the PHREEQC/ PHT3D database. We assume that the NAPL source
consists of the four BTEX compounds and that the remaining 
fraction can be represented by a fifth compound that is 
assumed insoluble. The computation of benzene 
dissolution within such a multicomponent NAPL mix can be programmed
as follows:
\\
{\color{red}
\begin{quote}
Benznapl \\
  -start  \\
\#  Compute kinetic dissolution from multicomp. NAPL: \\
   10  mBenznapl = tot("Benznapl")    \\
   15  if (mBenznapl $<=$ 1e-10) then goto 200  \\
   20  solub\_Benz = 0.022820  \\
   25  mBenz = tot("Benz")  \\
   32  mTolunapl = tot("Tolunapl")  \\
   34  mEthynapl = tot("Ethynapl")  \\
   36  mXylnapl  = tot("Xylnapl")  \\
   38  mLowsolubnapl  = tot("Lowsolubnapl")  \\
   40  m\_napl\_tot = mBenznapl + mTolunapl + mEthynapl \\
   41  m\_napl\_tot = m\_napl\_tot + mXylnapl + mLowsolubnapl  \\
   42  if (m\_napl\_tot $<=$ 1e-10) then goto 200  \\
   50  msolub\_Benz =  mBenznapl $/$ m\_napl\_tot $*$ solub\_Benz  \\
   60  rate = parm(2) $/$ 86400  $*$ (msolub\_Benz - mBenz)\\                                   
   70  moles = rate $*$ time  \\
   80  if (moles $>$ m) then moles = m  \\
   200 save moles\\
  -end  \\
 \end{quote}
 }
 
The above rate expression as well as those for the dissolution
of toluene, ethylbenzene and xylene have been included into the database
for this exercise. The values used for $C_{i}^{sat}$ are 
listed in Table \ref{solub_single_species}. Note the much higher
solubility of benzene, which is considered the most toxic of the
BTEX compounds.

\begin{table}[b]
\caption{Single species solubilities of BTEX compounds used for the simulations
in PHT3D Exercise 4.}
\begin{center}
\begin{tabular}{l c }
\toprule
Aqueous component    & $C_{i}^{sat}$     \\
                     & $(mol/l)$        \\ 
\midrule
Benzene              &  0.022820  \\
Toluene              &  0.005978  \\
Ethylbenzene         &  0.00193   \\
Xylene               &  0.001800  \\
\bottomrule

\end{tabular}
\label{solub_single_species}
\end{center}
\end{table}

To start the reactive transport modelling

\begin{itemize}

\item copy the two files {\color{red} \textbf{pht3d\_datab.ex\_eden\_hill}} (which contains the reaction database) and
{\color{red} \textbf{pmwin\_pht3dv2.ex\_eden\_hill}} into  
the database library folder {\color{red} \textbf{C:$\backslash$ \ldots $\backslash$pht3d$\backslash$database}} and include
the module in the file {\color{red} \textbf{pht3d\_module\_definition.txt}}
which lists all available reaction modules. 

\item Select the newly installed reaction module and define
a very simple water composition that contains only {\color{red} \textbf{Na}}, 
{\color{red} \textbf{Cl}}, {\color{red} \textbf{Benz}} (= dissolved benzene), 
{\color{red} \textbf{Tolu}} (= dissolved toluene), {\color{red} \textbf{Ethy}} (= dissolved ethylbenzene) and {\color{red} \textbf{Xyl}} (= dissolved xylene).
Include also the species {\color{red} \textbf{Benznapl}}, {\color{red} \textbf{Tolunapl}}, {\color{red} \textbf{Ethynapl}}, {\color{red} \textbf{Xylnapl}} and {\color{red} \textbf{Lowsolubnapl}} which together represent the (immobile) NAPL source.  

\item For both the ambient water and the recharge water define concentrations 
of 0.001 $mol/l$ for {\color{red} \textbf{Na}} and {\color{red} \textbf{Cl}}, respectively. 
Set the initial values (initial concentrations) of {\color{red} \textbf{pH}} to 7 
and set the {\color{red} \textbf{pe}} to 4.0. 

\item Define a NAPL-polluted source area near the water table (20 m $\geq$ z $\geq$ 19 m) 
between x = $15 m$ and x = $35 m$. Select an initial concentration of
0.018 $mol/l$ for the 4 NAPL compounds representing BTEX. Enter an initial concentration
of 0.072 $mol/l$ for the compound {\color{red} \textbf{Lowsolubnapl}} to all the grid cells
that sit above a level of 19 m.
Note, that the concentration of these immobile NAPL compounds
is defined in units of $mol/l$.

\item Define the settings for advective transport: Use the MMOC scheme in this and other 
PHT3D simulations that include free water tables. 

\item Define a longitudinal dispersivity of $0.05m$ and a transversal 
dispersivity of $5mm$.

\item Select Models $\rightarrow$ PHT3D $\rightarrow$ Simulation Settings 
$\rightarrow$ OK $\rightarrow$ Options $\rightarrow$ Threshold values for executing PHREEQC $\rightarrow$ and set the value for \textbf{Changes in pH} to 0.0001.

\item Run the simulation and visualise the results. 

\item Define some observation points and visualise the breakthrough curves
of benzene and xylene.

\item What explains the different characteristics of the BTCs ?

\end{itemize}

In the next part of this exercise we use instead of the Na-Cl solution
the measured and charge-balanced water composition from the Eden Hill site as given in \cite{prommer_jch} and as listed in Table \ref{eh_caqu}.
The water is anaerobic and only sulphate will act as electron acceptor. 
As discussed in the earlier PHREEQC example/exercise we will include 
microbial growth and decay in the simulations. Microbial growth
will only occur where both toluene and sulfate are present simultaneously.
The rate expression used in the PHREEQC Exercise 5 was slightly 
modified and compared to Eqn. (\ref{eq:v_m_ea}) an additional biomass 
inhibition term was included:

\begin{equation}
v_{m} = v_{max} I_{bio} \frac{C_{org}}{K_{org} + C_{org}} \frac{C_{ea}}{K_{ea} + C_{ea}} X
\end{equation}

This additional term $I_{bio}$ goes back to some conceptual models for biologically 
mediated degradation reactions which suggest that with increasing 
biomass concentrations and thus increasing biofilm thickness, 
degradation rates might become limited by the supply of reactants. 
In order to avoid solving the diffusional transport of reactants 
at a microscopic (i.e., within biofilm) level, a macroscopic formulation 
can be used to account for the rate-limitation resulting 
from excessive biomass accumulation. 
The biomass inhibition term is computed from

\begin{equation}
{ I_{bio} = \frac{\theta_{bio,max} - X_{tot}}{\theta_{bio,max}} }
\label{eq:lambda_diff}
\end{equation}
%

The above equation was used as a base for the computation 
of the degradation rates with PHREEQC/PHT3D: 

{\color{red}
\begin{quote}
Tolu \\
  -start \\
   10 mSulfred  = tot("Sulfred")  \\
   20 mSulf     = tot("S(6)")  \\
   30 mTolu     = tot("Tolu")  \\
   35 if (mTolu  $<$ 1e-10) then goto 220  \\
   40 if (mSulf  $<$ 1e-10) then goto 220  \\
   68 v\_up\_max\_Tolu = parm(1) $/$ 86400  \\
   84 theta\_bio\_max  = parm(2)  \\
   90 inhib\_bio = (theta\_bio\_max-mSulfred) $/$ theta\_bio\_max  \\
   95 monSulf = mSulf $/$ (0.0001 + mSulf)  \\
  100 monTolu = mTolu $/$ (0.0001 + mTolu)  \\
  150 v\_up\_Tolu = v\_up\_max\_Tolu $*$ monTolu $*$ monSulf  \\
  160 v\_up\_Tolu = v\_up\_Tolu $*$ inhib\_bio \\
  190 rate = v\_up\_Tolu $*$ mSulfred  \\
  200 moles = rate $*$ time  \\
  220 SAVE moles  \\
 -end  \\
\end{quote}
}

To proceed and include the biodegradation reactions in the simulations:

\begin{itemize}

\item Select Models $\rightarrow$ PHT3D $\rightarrow$ Simulation Settings and
activate the relevant components to simulate the water composition listed in
Table \ref{eh_caqu}. Activate also
sulfate reducing bacteria, i.e., {\color{red} \textbf{Sulfred}}.

\item Define the initial concentrations of the aqueous components
and the recharge water composition according to Table \ref{eh_caqu}.

\item Enter a value of 1e-07 $mol/l$ as initial concentration for 
the sulfate reducing bacteria. 

\item Check if the boundary conditions for the transport are defined correctly 
such that the inflowing water at the upstream boundary has the composition
of the ambient and the recharge water.  

\item Allocate the reaction rate parameters listed in 
Table \ref{tab:eh_rate_const}. To enter the values select
\item Select Models $\rightarrow$ PHT3D $\rightarrow$ Simulation Settings

\item Run the model

\item Visualise the $microbial$ $fringe$.

\end{itemize}

\begin{table}
\caption{Measured and equilibrated aqueous concentrations of 
the ambient water composition at the Eden Hill field site.}
\begin{center}
\begin{tabular}{l c }
\toprule
Aqueous component & $C_{init},C_{rech}$      \\
                  & $(mol\ l^{-1}_w)$        \\ 
\midrule
pH              &  5.14               \\
pe              &  8.01               \\
O(0)            &             0              \\
N$_{tot}$       &  1.00 $\times$ $10^{-4}$   \\ 
C(4)            &  7.96 $\times$ $10^{-3}$   \\
C(-4)            &            0              \\
S(6)            &  7.85 $\times$ $10^{-4}$   \\
S(-2)           &             0              \\
Fe(2)           &  5.20 $\times$ $10^{-5}$   \\
Fe(3)           &  1.46 $\times$ $10^{-5}$   \\
Ca              &  8.42 $\times$ $10^{-4}$   \\
K               &  2.59 $\times$ $10^{-4}$   \\
Mg              &  5.67 $\times$ $10^{-4}$   \\
Na              &  5.25 $\times$ $10^{-3}$   \\
Cl              &  6.46 $\times$ $10^{-3}$   \\
\bottomrule

\end{tabular}
\label{eh_caqu}
\end{center}
\end{table}

\begin{table}[hb]
\topcaption{Rate constants for the reactive transport simulation of BTEX degradation
under sulfate reducing conditions.}

\centering
\footnotesize
\begin{tabular*}{\textwidth}{@{\extracolsep{\fill}}lcccc}
\toprule
Rate Expression & Parameter Nr & Parameter & Unit & Value \\
\midrule
Benz        & 1 & $v_{max,sulf,benzene}$     & (1/day)     & 0.0   \\
Tolu        & 1 & $v_{max,sulf,toluene}$     & (1/day)     & 5.0   \\
Ethy        & 1 & $v_{max,sulf,ethylbenzene}$& (1/day)     & 2.0   \\
Xyl         & 1 & $v_{max,sulf,xylene}   $   & (1/day)     & 2.0   \\
Benz        & 2 & $\theta_{bio,max}$ & (mol/l)              &  $1.0 \cdot 10^{-5}$  \\
Tolu        & 2 & $\theta_{bio,max}$ & (mol/l)              &  $1.0 \cdot 10^{-5}$  \\
Ethy        & 2 & $\theta_{bio,max}$ & (mol/l)              &  $1.0 \cdot 10^{-5}$  \\
Xyl         & 2 & $\theta_{bio,max}$ & (mol/l)              &  $1.0 \cdot 10^{-5}$  \\
Benznapl    & 1 & not used         & -   &  1  \\
Tolunapl    & 1 & not used         & -   &  1  \\
Ethynapl    & 1 & not used         & -   &  1  \\
Xylnapl     & 1 & not used         & -   &  1  \\
Benznapl    & 2 & $\omega$  & (1/day)       & 0.025   \\
Tolunapl    & 2 & $\omega$  & (1/day)       & 0.025   \\
Ethynapl    & 2 & $\omega$  & (1/day)       & 0.025   \\
Xylnapl     & 2 & $\omega$  & (1/day)       & 0.025   \\
Sulfred & 1 & $v_{dec} $    & (1/day)                      &   0.1   \\
Sulfred & 2 & $C_{init}$    & (mol/l)                      &  $1.0 \cdot 10^{-7}$  \\
\bottomrule
\end{tabular*}
\label{tab:eh_rate_const}
\end{table}